\documentclass[12pt]{article}
\usepackage{amsmath}
\usepackage{geometry}
\geometry{margin=1.4in}
\usepackage{hyperref}
\usepackage{setspace}
\setstretch{1.4}

\title{Politics as Displacement Dynamics:\\
Governance, Power, and the Wrong Ground State}
\author{Diego Rincón\\
\small Independent}
\date{}

\begin{document}

\maketitle

\begin{abstract}
Political systems cycle through attractors: democracy becomes oligarchy, revolution becomes tyranny, idealism becomes pragmatism. Each transition is not corruption but displacement. Power finds a stable ground state that is cheaper to maintain than the one the system was designed for. The cost of returning to the intended ground state exceeds the cost of staying at the current attractor.
\end{abstract}

\section{Democracy as Expensive Ground State}

Democracy is the intended ground state: distributed power, accountability, participation. It is also expensive. Maintaining it requires constant friction (debate, opposition, resistance).

Power naturally flows to concentration $S^*$ because it is cheaper. A smaller set of decision-makers requires less coordination. Fewer voices means faster action.

The system does not fail. It returns to a lower-cost ground state: oligarchy, autocracy, or power captured by narrow interest.

\section{The Return Path}

Once a political system reaches $S^*$ (concentrated power), the return to $S_0$ (distributed power) is expensive. It requires:
- Dismantling existing power structures
- Overcoming inertia of the current attractor
- Building new institutions from scratch

The cost is so high that the system stays at $S^*$ unless forced by crisis.

\section{Revolutions as Failed Returns}

A revolution attempts to return to $S_0$ but is exhausting. Once the old ground state is destroyed, the system must find a new one. In the absence of careful design, it finds $S^{**}$: a different wrong attractor (often more oppressive than the original).

The pattern: revolution → tyranny → reform → oligarchy → stagnation → revolution.

Each is a different stable state. None is the intended ground.

\end{document}
